\chapter{المتجهات}\label{ch:g10:vectors}

تشفّر \omterm{def:g10:vectors:vector}{المتجهة} انتقالًا: اتجاهًا وطولًا، أيًّا كانت
نقطة الانطلاق. وتعطي المتجهات براهين نظيفة لوقائع هندسية، وتردّها،
متى دخلت \omterm{def:g10:coordgeom:system}{الإحداثيات} الصورة، إلى حسابات صغيرة. وستكون أداة أساسية
في \cref{ch:g10:lines}، ثم لاحقًا من أجل الجداء السلّمي
في \cref{ch:g11:scal}.

\section{الانسحابات والمتجهات}

\begin{definition}[المتجهة]\label{def:g10:vectors:vector}
إذا أُعطيت نقطتان $A$ و $B$، فإن \emph{المتجهة}\index{متجهة} $\vect{AB}$
تمثّل الانتقال من $A$ إلى $B$: ولها
\emph{اتجاه} (اتجاه المستقيم $(AB)$، مع جهة السير
عليه، من $A$ نحو $B$) و
\emph{طول} (المسافة $AB$، ويُكتب أيضًا $\norm{\vect{AB}}$).
وتكون متجهتان \emph{متساويتين} إذا شفّرتا الانتقال نفسه:
فالكتابة $\vect{AB} = \vect{CD}$ تعني أن تطبيق الانتقالين على أي
نقطة يعطي النتيجة نفسها.
\end{definition}

\begin{proposition}[المتجهات المتساوية ومتوازيات الأضلاع]\label{prop:g10:vectors:parallelogram}
تتحقق $\vect{AB} = \vect{CD}$ بالضبط عندما يكون $ABDC$ (بهذا \omterm{def:g10:coordgeom:system}{الترتيب}!)
متوازي أضلاع، وقد يكون مفلطحًا --- أي، بصورة مكافئة، عندما يكون للقطعتين $[AD]$
و $[BC]$ \omterm{prop:g10:coordgeom:midpoint}{المنتصف} نفسه.
\end{proposition}
\admitted

\begin{omfigure}
\begin{tikzpicture}[scale=1.0]
  \coordinate (A) at (0,0);
  \coordinate (B) at (2.2,0.8);
  \coordinate (C) at (1.0,-1.2);
  \coordinate (D) at (3.2,-0.4);
  \draw[omExo, thin] (A) -- (C);
  \draw[omExo, thin] (B) -- (D);
  \draw[-Stealth, omDef, thick] (A) -- (B)
    node[midway, above, font=\small] {$\vect{AB}$};
  \draw[-Stealth, omDef, thick] (C) -- (D)
    node[midway, below, font=\small] {$\vect{CD}$};
  \fill (A) circle (1.5pt) node[left, font=\small] {$A$};
  \fill (B) circle (1.5pt) node[right, font=\small] {$B$};
  \fill (C) circle (1.5pt) node[left, font=\small] {$C$};
  \fill (D) circle (1.5pt) node[right, font=\small] {$D$};
\end{tikzpicture}

{\small متجهتان متساويتان $\vect{AB} = \vect{CD}$: الاتجاه نفسه، والجهة نفسها،
والطول نفسه. والرباعي $ABDC$ متوازي أضلاع.}
\end{omfigure}

\begin{notation}
كثيرًا ما تُسمّى \omterm{def:g10:vectors:vector}{المتجهة} بحرف واحد، $\vec u$ أو $\vec v$، عندما لا
تهم النقطتان الطرفيتان. و\emph{\omterm{def:g10:vectors:vector}{المتجهة} المعدومة} $\vec 0 = \vect{AA}$ هي
الانتقال الذي لا يحرّك شيئًا. و\emph{مقابل} \omterm{def:g10:vectors:vector}{المتجهة}
$\vect{AB}$ هو $\vect{BA}$: الطول نفسه والجهة المعاكسة. ونكتب
$\vect{BA} = -\vect{AB}$.
\end{notation}

\section{جمع المتجهات}

\begin{definition}[مجموع متجهتين]\label{def:g10:vectors:sum}
\emph{المجموع}\index{متجهة!مجموع} $\vec u + \vec v$ هو الانتقال
الناتج عن إنجاز $\vec u$ ثم $\vec v$.
\end{definition}

\begin{theorem}[علاقة شال]\label{thm:g10:vectors:chasles}
من أجل أي ثلاث نقط $A$ و $B$ و $C$:
\[
\vect{AB} + \vect{BC} = \vect{AC}.
\]
\end{theorem}

\begin{proof}
الانتقال من $A$ إلى $B$، ثم من $B$ إلى $C$، انتقال ينقل
$A$ إلى $C$: أي أنه الانتقال الذي تشفّره $\vect{AC}$.
\end{proof}

\begin{remark}[قاعدة متوازي الأضلاع]
لجمع متجهتين مرسومتين من النقطة نفسها، $\vect{AB} + \vect{AC}$،
أتمّ متوازي الأضلاع $ABDC$: \omterm{def:g10:vectors:sum}{فالمجموع} هو القطر $\vect{AD}$.
وفعلًا $\vect{AC} = \vect{BD}$، إذن
$\vect{AB} + \vect{AC} = \vect{AB} + \vect{BD} = \vect{AD}$ حسب علاقة
شال.
\end{remark}

\begin{omfigure}
\begin{tikzpicture}[scale=1.05]
  \coordinate (A) at (0,0);
  \coordinate (B) at (2.4,0.5);
  \coordinate (C) at (3.4,2.0);
  \draw[-Stealth, omDef, thick] (A) -- (B)
    node[midway, below, font=\small] {$\vec u$};
  \draw[-Stealth, omThm, thick] (B) -- (C)
    node[midway, right, font=\small] {$\vec v$};
  \draw[-Stealth, omMeth, thick] (A) -- (C)
    node[midway, above left, font=\small] {$\vec u + \vec v$};
\end{tikzpicture}\hspace{1.4cm}%
\begin{tikzpicture}[scale=1.05]
  \coordinate (A) at (0,0);
  \coordinate (B) at (2.4,0.5);
  \coordinate (C) at (1.0,1.5);
  \coordinate (D) at (3.4,2.0);
  \draw[omExo, thin] (B) -- (D) -- (C);
  \draw[-Stealth, omDef, thick] (A) -- (B)
    node[midway, below, font=\small] {$\vect{AB}$};
  \draw[-Stealth, omThm, thick] (A) -- (C)
    node[midway, left, font=\small] {$\vect{AC}$};
  \draw[-Stealth, omMeth, thick] (A) -- (D)
    node[pos=0.75, below right, font=\small] {$\vect{AD}$};
  \fill (D) circle (1.2pt) node[right, font=\small] {$D$};
\end{tikzpicture}

{\small رسمان للمجموع: رأسًا إلى ذيل (علاقة شال، إلى اليسار) و
قاعدة متوازي الأضلاع (إلى اليمين).}
\end{omfigure}

\begin{example}[التبسيط بعلاقة شال]\label{ex:g10:vectors:chasles}
بسّط $\vect{MN} + \vect{NP} + \vect{PQ}$: تسلسل الانتقالات
يعطي $\vect{MQ}$. وبسّط $\vect{AB} - \vect{AC}$: أعد كتابة
الفرق مجموعًا مع \omterm{def:g10:vectors:vector}{المتجهة} المقابلة،
\[
\vect{AB} - \vect{AC} = \vect{AB} + \vect{CA}
= \vect{CA} + \vect{AB} = \vect{CB}.
\]
\end{example}

\section{ضرب متجهة في عدد}

\begin{definition}[الضرب في عدد]\label{def:g10:vectors:scalar}
لتكن $\vec u$ \omterm{def:g10:vectors:vector}{متجهة} غير معدومة وليكن $k$ عددًا حقيقيًا. \omterm{def:g10:vectors:vector}{للمتجهة}
$k\vec u$\index{متجهة!ضرب في عدد} اتجاه $\vec u$ نفسه،
وطولها $\abs{k} \times \norm{\vec u}$، وجهتها جهة
$\vec u$ إذا كان $k > 0$، والجهة المعاكسة إذا كان $k < 0$. ومن أجل $k = 0$،
$0\vec u = \vec 0$.
\end{definition}

\begin{definition}[الاستقامية]\label{def:g10:vectors:collinear}
تكون المتجهتان $\vec u$ و $\vec v$
\emph{مستقيميتين}\index{متجهتان مستقيميتان} إذا كان لهما الاتجاه
نفسه، أي إذا كانت $\vec v = k \vec u$ من أجل \omterm{def:g10:numbers:sets}{عدد حقيقي} $k$ ما (أو كانت
إحداهما $\vec 0$).
\end{definition}

\begin{remark}
الاستقامية هي لغة المتجهات للتوازي والاستقامة:
\begin{itemize}
  \item يكون المستقيمان $(AB)$ و $(CD)$ متوازيين بالضبط عندما تكون $\vect{AB}$
        و $\vect{CD}$ \omterm{def:g10:vectors:collinear}{مستقيميتين}؛
  \item وتكون النقط $A$ و $B$ و $C$ على استقامة واحدة بالضبط عندما تكون $\vect{AB}$ و
        $\vect{AC}$ \omterm{def:g10:vectors:collinear}{مستقيميتين}.
\end{itemize}
\end{remark}

\section{إحداثيات المتجهات}

\begin{definition}[إحداثيات متجهة]\label{def:g10:vectors:coords}
في \omterm{def:g10:coordgeom:system}{نظام إحداثيات}، \emph{\omterm{def:g10:coordgeom:system}{إحداثيات}} \omterm{def:g10:vectors:vector}{المتجهة} $\vect{AB}$
هي العددان اللذان يصفان الانتقال:
\[
\vect{AB}\,(x_B - x_A,\ y_B - y_A).
\]
\omterm{def:g10:vectors:vector}{فالمتجهة} $\vec u\,(a, b)$ تنقل كل نقطة $a$ وحدة أفقيًا و
$b$ وحدة شاقوليًا.
\end{definition}

\begin{proposition}[الحساب بالإحداثيات]\label{prop:g10:vectors:coordrules}
لتكن $\vec u\,(a, b)$ و $\vec v\,(c, d)$، وليكن $k \in \R$.
\begin{enumerate}
  \item تتحقق $\vec u = \vec v$ بالضبط عندما $a = c$ و $b = d$؛
  \item \omterm{def:g10:coordgeom:system}{إحداثيات} $\vec u + \vec v$ هي $(a + c,\ b + d)$؛
  \item \omterm{def:g10:coordgeom:system}{إحداثيات} $k\vec u$ هي $(ka,\ kb)$؛
  \item $\norm{\vec u} = \sqrt{a^2 + b^2}$ (في \omterm{def:g10:coordgeom:system}{نظام متعامد ومتجانس}).
\end{enumerate}
\end{proposition}

\begin{proof}
تعيد النقاط 1--3 صياغة ما تفعله الانتقالات، إحداثيًا إحداثيًا: فمثلًا
إنجاز $\vec u$ ثم $\vec v$ ينقل نقطة أفقيًا بمقدار
$a$ ثم $c$، أي بمقدار $a + c$ إجمالًا. أما النقطة 4 فهي صيغة المسافة
في \cref{thm:g10:coordgeom:distance} مطبقة على قطعة تمثّل
$\vec u$.
\end{proof}

\begin{theorem}[محك الاستقامية]\label{thm:g10:vectors:det}
تكون المتجهتان $\vec u\,(a, b)$ و $\vec v\,(c, d)$ \omterm{def:g10:vectors:collinear}{مستقيميتين} إذا
وفقط إذا
\[
ad - bc = 0 .
\]
\end{theorem}

\begin{proof}
إذا كانت $\vec v = k\vec u$، فإن $c = ka$ و $d = kb$، إذن
$ad - bc = a(kb) - b(ka) = 0$. ويصح الشيء نفسه إذا كانت $\vec u = k \vec v$ أو إذا كانت
إحدى المتجهتين معدومة.

وبالعكس، لنفترض أن $ad - bc = 0$ مع $\vec u \neq \vec 0$، وليكن مثلًا
$a \neq 0$ (والحالة $b \neq 0$ مماثلة). ضع $k = \frac{c}{a}$؛ عندئذٍ
$c = ka$، وتعطي $ad = bc = b(ka)$، بعد القسمة على $a \neq 0$،
أن $d = kb$. إذن $\vec v = k\vec u$.
\end{proof}

\begin{example}\label{ex:g10:vectors:det}
هل المتجهتان $\vec u\,(3, -2)$ و $\vec v\,(-6, 4)$ مستقيميتان؟ احسب
$ad - bc = 3 \times 4 - (-2)\times(-6) = 12 - 12 = 0$: نعم، وفعلًا
$\vec v = -2\vec u$. أما من أجل $\vec u\,(3, -2)$ و $\vec w\,(1, 5)$:
$3 \times 5 - (-2) \times 1 = 17 \neq 0$: فليستا \omterm{def:g10:vectors:collinear}{مستقيميتين}.
\end{example}

\begin{method}[الاستقامة والتوازي]\label{met:g10:vectors:alignment}
للبرهان على أن ثلاث نقط $A$ و $B$ و $C$ على استقامة واحدة:
\begin{enumerate}
  \item احسب \omterm{def:g10:coordgeom:system}{إحداثيات} $\vect{AB}$ و $\vect{AC}$؛
  \item تحقق من المحك $ad - bc = 0$ في \cref{thm:g10:vectors:det}؛
  \item استنتج: \omterm{def:g10:vectors:collinear}{المتجهتان مستقيميتان} وتشتركان في النقطة $A$، إذن
        النقط الثلاث على استقامة واحدة.
\end{enumerate}
والحساب نفسه مع $\vect{AB}$ و $\vect{CD}$ يبرهن على أن
المستقيمين $(AB)$ و $(CD)$ متوازيان.
\end{method}

\begin{example}\label{ex:g10:vectors:alignment}
لتكن $A(1, 2)$ و $B(3, 5)$ و $C(7, 11)$. عندئذٍ $\vect{AB}\,(2, 3)$ و
$\vect{AC}\,(6, 9)$، و $2 \times 9 - 3 \times 6 = 18 - 18 = 0$: إذن
النقط على استقامة واحدة (وفي الواقع $\vect{AC} = 3\vect{AB}$).
\end{example}

\begin{omfigure}
\begin{tikzpicture}[scale=0.62]
  \draw[thin, gray!40] (-0.4,-0.4) grid (8.4,5.4);
  \draw[-Stealth] (-0.6,0) -- (8.6,0);
  \draw[-Stealth] (0,-0.6) -- (0,5.6);
  \coordinate (A) at (1,1);
  \coordinate (B) at (3,2);
  \coordinate (C) at (7,4);
  \draw[-Stealth, omDef, very thick] (A) -- (B)
    node[midway, above left=-2pt, font=\small] {$\vect{AB}$};
  \draw[-Stealth, omThm, thick] (A) -- (C)
    node[pos=0.8, below right=-2pt, font=\small] {$\vect{AC} = 3\vect{AB}$};
  \fill (A) circle (2pt) node[above left, font=\small] {$A$};
  \fill (B) circle (2pt) node[above left, font=\small] {$B$};
  \fill (C) circle (2pt) node[above left, font=\small] {$C$};
\end{tikzpicture}

{\small نقط على استقامة واحدة: \omterm{def:g10:vectors:vector}{المتجهة} $\vect{AC}$ مضاعف سلّمي \omterm{def:g10:vectors:vector}{للمتجهة} $\vect{AB}$.}
\end{omfigure}

\section{تمارين}

\begin{exercise}[$\star$]\label{exo:g10:vectors:1}
بسّط مستعملًا علاقة شال:
\[
\vect{AB} + \vect{BD}, \qquad
\vect{KL} + \vect{LM} + \vect{MK}, \qquad
\vect{AC} - \vect{BC}, \qquad
\vect{AB} + \vect{CA}.
\]
\end{exercise}

\begin{exercise}[$\star$]\label{exo:g10:vectors:2}
لتكن $A(2, 1)$ و $B(5, 3)$ و $C(-1, 4)$. احسب \omterm{def:g10:coordgeom:system}{إحداثيات}
$\vect{AB}$ و $\vect{BC}$ و $\vect{AC}$، وتحقق على \omterm{def:g10:coordgeom:system}{الإحداثيات} من أن
$\vect{AB} + \vect{BC} = \vect{AC}$.
\end{exercise}

\begin{exercise}[$\star$]\label{exo:g10:vectors:3}
لتكن $\vec u\,(2, -3)$ و $\vec v\,(-1, 4)$. أعطِ \omterm{def:g10:coordgeom:system}{إحداثيات}
$\vec u + \vec v$ و $3\vec u$ و $2\vec u - \vec v$، واحسب
$\norm{\vec u}$.
\end{exercise}

\begin{exercise}[$\star$]\label{exo:g10:vectors:4}
في كل حالة، قل هل $\vec u$ و $\vec v$ مستقيميتان:
\[
\text{(أ) } \vec u\,(4, 6),\ \vec v\,(6, 9);
\qquad
\text{(ب) } \vec u\,(2, -5),\ \vec v\,(-4, 10);
\qquad
\text{(ج) } \vec u\,(3, 1),\ \vec v\,(1, 3).
\]
\end{exercise}

\begin{exercise}[$\star$]\label{exo:g10:vectors:5}
لتكن $A(0, 2)$ و $B(4, 0)$ و $C(6, 3)$. أوجد \omterm{def:g10:coordgeom:system}{إحداثيات} النقطة
$D$ التي تجعل $ABCD$ متوازي أضلاع، أي التي تحقق
$\vect{AB} = \vect{DC}$.
\end{exercise}

\begin{exercise}[$\star\star$]\label{exo:g10:vectors:6}
لتكن $A(-2, 1)$ و $B(1, 3)$ و $C(4, -1)$.
\begin{enumerate}
  \item احسب \omterm{def:g10:coordgeom:system}{إحداثيات} النقطة $M$ التي تحقق
        $\vect{AM} = \vect{AB} + \vect{AC}$.
  \item ما نوع الرباعي $ABMC$؟ برّر.
\end{enumerate}
\end{exercise}

\begin{exercise}[$\star\star$]\label{exo:g10:vectors:7}
لتكن $A(1, -1)$ و $B(4, 1)$ و $C(10, 5)$. هل النقط $A$ و $B$ و $C$
على استقامة واحدة؟ والسؤال نفسه من أجل $A(0, 3)$ و $B(2, 2)$ و $C(8, -1)$.
\end{exercise}

\begin{exercise}[$\star\star$]\label{exo:g10:vectors:8}
لتكن $A(1, 2)$ و $B(5, 4)$ و $C(6, 1)$ و $D(2, -1)$. بيّن أن $(AB)$ و
$(CD)$ متوازيان، ثم أن $ABCD$ متوازي أضلاع. وأي
تحقق يستلزم الآخر؟
\end{exercise}

\begin{exercise}[$\star\star$]\label{exo:g10:vectors:9}
ليكن $ABC$ مثلثًا، ولتكن $I$ النقطة المعرَّفة بالعلاقة
$\vect{AI} = \frac23 \vect{AB}$. مع $A(0,0)$ و $B(6, 3)$ و $C(2, 5)$
(حتى تبقى الحسابات بسيطة):
\begin{enumerate}
  \item احسب \omterm{def:g10:coordgeom:system}{إحداثيات} $I$؛
  \item لتكن $J$ بحيث $\vect{CJ} = \frac23\vect{CB}$؛ احسب
        \omterm{def:g10:coordgeom:system}{إحداثيات} $J$؛
  \item بيّن أن $(IJ)$ يوازي $(AC)$.
\end{enumerate}
\end{exercise}

\begin{exercise}[$\star\star\star$]\label{exo:g10:vectors:10}
ليكن $ABCD$ رباعيًا كيفيًا، ولتكن $I$ و $J$ و $K$ و $L$
منتصفات $[AB]$ و $[BC]$ و $[CD]$ و $[DA]$. مستعملًا \omterm{def:g10:coordgeom:system}{الإحداثيات}
$A(x_A, y_A)$ وغيرها، بيّن أن $\vect{IJ} = \vect{LK}$، واستنتج أن
منتصفات أضلاع \emph{أي} رباعي تشكّل
متوازي أضلاع.
\end{exercise}

\section{مسألة: جبر الأسهم}

\begin{problem}[{مسألة نهاية الأسبوع --- تمارين شال، ومركز
الثقل معادًا برهانه في سطرين، والقوى في توازن: ما
المتجهات في الحقيقة}]
\label{pb:g10:vectors:1}
برهن الكتاب السابق على أن متوسطات المثلث تلتقي عند ثلثي
طولها --- بمتوازي أضلاع بارع مخبّأ داخل
المثلث (مسألة نهاية الأسبوع عن مركز الثقل في الكتاب السابق). وتعيد المتجهات برهانه في
سطرين، وتعطي التلاقي مجانًا. وذلك ما يخدمه
هذا الجبر الجديد: فالمبرهنات تصير حسابات بالأسهم.
تمرّن هذه المسألة على الحساب (وعلى شال قبل كل شيء)،
وتقدّم البرهان ذا السطرين، وتوسّع فارينيون
(\cref{exo:g10:vectors:10})، وتنتهي حيث وُلدت المتجهات:
عند القوى والسرعات.

\medskip
\noindent\textbf{الجزء الأول --- تمارين شال.}
\begin{enumerate}
  \item بسّط مستعملًا علاقة شال
        (\cref{thm:g10:vectors:chasles}):
        $\vect{AB} + \vect{BC} + \vect{CD}$؛ \quad
        $\vect{MA} - \vect{MB}$؛ \quad
        $\vect{AB} + \vect{CD} + \vect{BC} + \vect{DA}$.
  \item حيلة المبدأ: بيّن أنه من أجل \emph{أي} نقطة $O$،
        $\vect{AB} = \vect{OB} - \vect{OA}$.
  \item برهن على التوصيفين التاليين \omterm{prop:g10:coordgeom:midpoint}{لمنتصف} $[AB]$،
        وليكن $I$:
        \[
        \vect{IA} + \vect{IB} = \vec 0
        \qquad\text{و}\qquad
        \vect{OI} = \tfrac12\left(\vect{OA} + \vect{OB}\right)
        \ \text{من أجل أي } O .
        \]
  \item مستعملًا $\vect{AB} = \vect{DC}$، أعد إيجاد الرأس الرابع
        $D$ لمتوازي الأضلاع $ABCD$ حيث $A(-1, 2)$ و
        $B(3, 4)$ و $C(6, 0)$ --- وقارن بجواب
        حيلة \omterm{prop:g10:coordgeom:midpoint}{المنتصف} في \cref{pb:g10:coordgeom:1}.
  \item إضافة \omterm{def:g10:vectors:vector}{متجهة} ثابتة $\vec u$ إلى كل نقطة من
        المستوي تنجز أي تحويل --- ذاك الذي
        اكتُشف بمرآتين في مسألة نهاية الأسبوع عن المرآتين في الكتاب السابق
        وبنصفي دورة في مسألة نهاية الأسبوع عن أنصاف الدورات في الكتاب السابق؟ احسب
        صورة $(1, 2)$ بالتحويل $\vec u = (3, -1)$.
\end{enumerate}

\medskip
\noindent\textbf{الجزء الثاني --- مركز الثقل، برهان ثالث.}
نعرّف النقطة $G$ من المثلث $ABC$ بالشرط
الأنيق
\[
\vect{GA} + \vect{GB} + \vect{GC} = \vec 0 .
\]
\begin{enumerate}[resume]
  \item مستعملًا حيلة المبدأ، بيّن أن هذا الشرط
        يحدد $G$ تحديدًا وحيدًا:
        $\vect{OG} = \frac13\left(\vect{OA} + \vect{OB} +
        \vect{OC}\right)$ من أجل أي $O$ --- إذن \omterm{def:g10:coordgeom:system}{إحداثيات}
        $G$ هي \emph{متوسطات} \omterm{def:g10:coordgeom:system}{إحداثيات} الرؤوس
        (وقد رأت ذلك مسألة نهاية الأسبوع عن مركز الثقل في الكتاب السابق عدديًا).
  \item برهن على $\vect{AG} = \frac13\left(\vect{AB} +
        \vect{AC}\right)$، ثم، بكتابة $K$ \omterm{prop:g10:coordgeom:midpoint}{لمنتصف}
        $[BC]$، على $\vect{AK} = \frac12\left(\vect{AB} +
        \vect{AC}\right)$. واستنتج:
        $\vect{AG} = \frac23\,\vect{AK}$ --- أي مبرهنة
        الثلثين، معادًا برهانها.
  \item لماذا يمر المتوسطان من $B$ ومن $C$ بالنقطة
        $G$ \emph{نفسها} دون أي حساب إضافي؟
        قارن بين البراهين الثلاثة المتاحة لهذه المبرهنة ---
        متوازي الأضلاع المخبّأ في الكتاب السابق، \omterm{def:g10:coordgeom:system}{والإحداثيات}، والمتجهات ---
        في جملة واحدة لكل منها.
  \item تحقق عددي: $A(1, 1)$ و $B(5, 2)$ و $C(3, 6)$.
        احسب $G$، ثم $K$، ثم تحقق من
        $\vect{AG} = \frac23 \vect{AK}$ بمحك الاستقامية
        (\cref{thm:g10:vectors:det}).
  \item تقرأ الفيزياء $G$ نقطةَ توازن لثلاث كتل
        متساوية عند الرؤوس. ضاعف الكتلة عند $A$ (الكتل
        $2, 1, 1$): فتصير نقطة التوازن
        $\vect{OG'} = \frac{2\vect{OA} + \vect{OB} +
        \vect{OC}}{4}$. احسب $G'$ من أجل مثلث
        السؤال 9 وقارن موضعه بموضع $G$.
\end{enumerate}

\medskip
\noindent\textbf{الجزء الثالث --- الاستقامية في العمل.}
\begin{enumerate}[resume]
  \item هل $\vec u(3, -2)$ و $\vec v(-6, 4)$ مستقيميتان؟ وماذا عن
        $\vec w(2, 5)$ و $\vec z(4, 9)$؟ قرّر بمحك
        المحدد.
  \item هل النقط $A(1, 2)$ و $B(3, 5)$ و $C(7, 11)$
        على استقامة واحدة؟
  \item ما بعد فارينيون (\cref{exo:g10:vectors:10}): في أي
        رباعي $ABCD$، يصل \emph{المتوسطان الثنائيان} بين
        منتصفات الأضلاع المتقابلة ($[IK]$ و $[JL]$). مستعملًا
        متجهات الموضع ($\vect{OI} = \frac12(\vect{OA} +
        \vect{OB})$، وغيرها)، بيّن أن للمتوسطين الثنائيين
        \omterm{prop:g10:coordgeom:midpoint}{المنتصف} نفسه، \omterm{def:g10:vectors:vector}{ومتجهة} موضعه
        $\frac14\left(\vect{OA} + \vect{OB} + \vect{OC} +
        \vect{OD}\right)$: أي مركز ثقل الرباعي.
  \item طاليس في سطر واحد: $M$ و $N$ هما النقطتان المعرَّفتان
        بالعلاقتين $\vect{AM} = \frac13\vect{AB}$ و
        $\vect{AN} = \frac13\vect{AC}$. احسب
        $\vect{MN}$ بدلالة $\vect{BC}$ واستنتج
        التوازي والنسبة.
  \item عمّم السؤال 14 على نسبة كيفية $k$، وبيّن
        في جملة واحدة كيف يستوعب حساب المتجهات
        مبرهنة المنتصفين ومبرهنة
        طاليس دفعة واحدة.
\end{enumerate}

\medskip
\noindent\textbf{الجزء الرابع --- القوى والسرعات.}
\begin{enumerate}[resume]
  \item يدفع تيار قاربًا بسرعة $3$ كم/سا نحو الشرق
        بينما يجدّف بسرعة $4$ كم/سا نحو الشمال. أعطِ \omterm{def:g10:vectors:vector}{متجهة}
        السرعة المحصلة وطويلتها، وصف المسار
        الفعلي. (يظهر ثلاثي قديم جدًّا.)
  \item تؤثر ثلاث قوى في حلقة: $\vec u(2, 1)$ و
        $\vec v(-3, 2)$ و $\vec w(1, -3)$. احسب المحصلة.
        ماذا تفعل الحلقة؟ ولماذا يكون مجموع الأسهم
        على أي مضلع \emph{مغلق} مساويًا دائمًا $\vec 0$
        (شال)؟
  \item تؤثر قوتان $\vec u(4, -1)$ و $\vec v(-1, 3)$ في
        خطّاف. ما القوة الثالثة التي تبقيه في توازن؟
  \item يعبر سابح نهرًا متجهًا مباشرة إلى الضفة المقابلة بسرعة
        $2$ م/ثا؛ ويجري التيار بسرعة $1.5$ م/ثا. أعطِ
        السرعة المحصلة وطويلتها. وعرض النهر
        $50$ م: كم يستغرق العبور، وكم
        ينزاح السابح مع التيار؟
  \item الخاتمة: كائن واحد في ثلاثة أثواب --- سهم
        انتقال، وزوج \omterm{def:g10:coordgeom:system}{إحداثيات}، وقوة أو سرعة. اذكر
        المبرهنات الكلاسيكية الثلاث التي أعادت هذه المسألة برهانها
        بجبر المتجهات، وسمِّ المقدار الهندسي الوحيد الذي لا
        تستطيع متجهات هذا الفصل قياسه --- فالأداة التي
        تعالج ذلك تأتي مع الجداء السلّمي، في السنة المقبلة.
\end{enumerate}
\end{problem}
